\section{Objetivos}

\subsection{Objetivo general}

Desarrollar un prototipo de escáner LIDAR 2D de bajo costo que construya mapas
de su entorno aplicando transformaciones lineales ---rotaciones en
$\mathbb{R}^2$ y transformaciones afines en coordenadas homogéneas--- y que
muestre en tiempo real los cálculos matriciales involucrados.

\subsection{Objetivos específicos}

\begin{enumerate}
\item Construir el prototipo físico con un ESP32, un sensor VL53L0X y un
servomotor SG90, junto con su firmware, para obtener lecturas $(r,\theta)$ en
un barrido de $0^\circ$ a $180^\circ$.

\item Implementar la conversión de coordenadas polares a cartesianas mediante
la matriz de rotación y la fusión de escaneos de dos posiciones mediante una
matriz homogénea de $3\times3$, verificando numéricamente sus propiedades:
ortogonalidad, determinante unitario, conservación de la norma y existencia de
la inversa.

\item Desarrollar una interfaz web que dibuje el mapa en tiempo real y muestre
paso a paso la multiplicación de matrices de cada lectura.

\item Cuantificar el efecto del ruido del sensor y de los errores en la
posición y orientación del aparato sobre la precisión del mapa fusionado.
\end{enumerate}
